\documentclass[10pt,a4paper]{article}
\usepackage{amsmath,amssymb,amsthm,graphicx}
\usepackage{hyperref}

\usepackage[scr=rsfs,cal=euler]{mathalfa}
\usepackage{bbm}
\let\mathds\mathbbm
\multlinegap0pt

\newenvironment{enumeratei}
{\bgroup\def\theenumi{\roman{enumi}}\def\theenumii{\arabic{enumii}}\begin{enumerate}}
{\end{enumerate}\egroup}
\newenvironment{enumerateI}
{\bgroup\def\theenumi{\Roman{enumi}}\def\theenumii{\arabic{enumii}}\begin{enumerate}}
{\end{enumerate}\egroup}

\newcommand{\Psfrac}[2]{(\sfrac{#1}{#2})}
\newcommand{\psfrac}[2]{\sfrac{(#1)}{#2}}
\newcommand{\spfrac}[2]{\sfrac{#1}{(#2)}}
\newcommand{\pspfrac}[2]{\sfrac{(#1)}{(#2)}}
\newcommand\mto{\mathchoice{\longmapsto}{\mapsto}{\mapsto}{\mapsto}}

\usepackage{mathtools}
\usepackage{esint}

\theoremstyle{plain}
\newtheorem{theor10}{Theorem}
\newenvironment{theor1}
{\renewcommand{\qedsymbol}{$\lozenge$}\pushQED{\qed}\begin{theor10}}
{\popQED\end{theor10}}
\newtheorem{cor10}{Corollary}
\newenvironment{cor1}
{\renewcommand{\qedsymbol}{$\lozenge$}\pushQED{\qed}\begin{cor10}}
{\popQED\end{cor10}}
\newtheorem{theor0}{Theorem}[section]
\newenvironment{theor}
{\renewcommand{\qedsymbol}{$\lozenge$}\pushQED{\qed}\begin{theor0}}
{\popQED\end{theor0}}
\newtheorem{lem0}[theor0]{Lemma}
\newenvironment{lem}
{\renewcommand{\qedsymbol}{$\lozenge$}\pushQED{\qed}\begin{lem0}}
{\popQED\end{lem0}}
\newtheorem{prop0}[theor0]{Proposition}
\newenvironment{prop}
{\renewcommand{\qedsymbol}{$\lozenge$}\pushQED{\qed}\begin{prop0}}
{\popQED\end{prop0}}
\newtheorem{cor0}[theor0]{Corollary}
\newenvironment{cor}
{\renewcommand{\qedsymbol}{$\lozenge$}\pushQED{\qed}\begin{cor0}}
{\popQED\end{cor0}}
\theoremstyle{definition}
\newtheorem{defin0}[theor0]{Definition}
\newenvironment{defin}
{\renewcommand{\qedsymbol}{$\lozenge$}\pushQED{\qed}\begin{defin0}}
{\popQED\end{defin0}}
\newtheorem{rems0}[theor0]{Remarks}
\newenvironment{rems}
{\renewcommand{\qedsymbol}{$\lozenge$}\pushQED{\qed}\begin{rems0}}
{\popQED\end{rems0}}
\newtheorem{rem0}[theor0]{Remark}
\newenvironment{rem}
{\renewcommand{\qedsymbol}{$\lozenge$}\pushQED{\qed}\begin{rem0}}
{\popQED\end{rem0}}

\numberwithin{equation}{section}

\newcommand{\N}{\mathbb N}
\newcommand{\e}{\varepsilon}
\newcommand{\Lc}{\mathcal{L}}
\newcommand{\Vc}{\mathcal V}
\newcommand{\R}{\mathbb R}
\newcommand{\E}{\mathbb E}
\newcommand{\C}{\mathbb C}
\newcommand{\F}{\mathcal F}
\newcommand{\loc}{{\operatorname{loc}}}
\newcommand{\Id}{\operatorname{Id}}
\newcommand{\per}{{\operatorname{per}}}
\newcommand{\ee}{e}
\newcommand{\Aa}{\boldsymbol a}
\newcommand{\bb}{{\boldsymbol b}}
\newcommand{\Bb}{{\boldsymbol B}}
\newcommand{\aA}{{\boldsymbol A}}
\newcommand{\cc}{{\boldsymbol c}}
\newcommand{\Ld}{\operatorname{L}}
\newcommand{\supp}{\operatorname{supp}}
\newcommand{\diam}{\operatorname{diam}}
\newcommand{\step}[1]{\subsubsection*{Step \textup #1}}
\newcommand{\Step}[1]{\Subsubsection*{Step \textup #1}}
\newcommand{\Pm}{\mathbb{P}}
\newcommand{\expec}[1]{\mathbb{E}\left[ #1 \right]}
\newcommand{\Bexpec}[1]{\mathbb{E}\Bigl[ #1 \Bigr]}
\newcommand{\expecm}[1]{\mathbb{E}\bigl[ #1 \bigr]}\let\bexpec\expecm

\usepackage{cleveref}
\crefname{originalequation}{original equation}{original equations}
\Crefname{originalequation}{Original equation}{Original equations}
\makeatletter
\newlabel{e.spec-1}{{1.1}{source}{}{sourcepointer-1}{}}
\newlabel{e.spec-1@cref}{{[originalequation][0][]1.1}{[1][0][]source}{}{}{}}
\makeatother
\title{The recursively defined corrector fields are well posed by the choice of their integration constant}
\author{Mitia Duerinckx, Antoine Gloria, Matthias Ruf}
\date{doi:10.5802/jep.259; CC BY 4.0}
\begin{document}
\maketitle
\setcounter{section}{1}
\noindent\textit{Editorial note.} All delivered mathematical text is the author's own, copied from the frozen source sources/jep-2024-259-a-spectral-ansatz-for-the-long-time-homogenization/pinned-original-source.tex; no statement or proof step was rewritten, reordered or removed.\par\medskip
\noindent\textit{Editorial note.} The unit has no row in the census, so the gate is the only authority for this card, printed numbers are computed by hand and the finished PDF is verified with pdftotext (duty desk rule for this band).\par\medskip
\noindent\textit{Editorial note.} The counted span carries no \textbackslash{}cite.\par\medskip
\noindent\textit{Editorial scope.} Single Lemma card cut from doi:10.5802/jep.259 (A spectral ansatz for the long-time homogenization of the wave equation); statement at source lines 830-848 and the author's own proof at source lines 849-878.\par\medskip
\noindent\textit{Allowed standard prerequisites (editorial).}\begin{itemize}
\item The recursively defined fields and the defining equation of the correctors. Same article, the recursion and the equation labelled e.zeta0n used inside the statement (source lines 830-836). It fixes the objects whose well-posedness the lemma asserts, together with the equation that the proof inserts into the averaged identity.
\item The defining equation of the auxiliary fields. Same article, the equation labelled e.spec-1 tested in the first display of the proof (source line 852). It is tested against the corrector to produce the averaged identity that drives the induction.
\item The Fredholm alternative on the periodic cell. The Fredholm alternative invoked in the statement at line 840. It is what turns vanishing average of the right hand side into well-posedness of the periodic problem.
\end{itemize}\medskip
\ifdefined\hypertarget\hypertarget{sourcepointer-1}{}\fi
\noindent\textit{Source reference.} Original equation 1.1: sources/jep-2024-259-a-spectral-ansatz-for-the-long-time-homogenization/pinned-original-source.tex, source lines 761--761. navigation stub: the target is the defining equation of the spectral correctors, outside the counted spans. Printed number computed by hand because the unit has no census row: it is the first numbered equation environment in the pinned source (line 761) and no section marker precedes it, so it is numbered within the article's first scope as (1.1).\par\medskip
\begin{lem}\label{lem:spec-corr-zeta-1}
For all $n\ge0$ and $\xi\in \R^d$, we recursively define $\check\zeta^{n,0}_\xi\in H^1_{{\per}}(Q)$ as the unique periodic scalar field that satisfies
\begin{multline}\label{e.zeta0n}
-\nabla\cdot\Aa\nabla\check \zeta^{n,0}_\xi=\nabla\cdot(\Aa i\xi\check\zeta^{n-1,0}_\xi)
+i\xi\cdot\Aa\big(\nabla\check\zeta^{n-1,0}_\xi+i\xi\check\zeta^{n-2,0}_\xi\big)\\
-\check\psi^{n+1}_\xi+\sum_{k=0}^{n+1}\Bexpec{\overline{\check\psi^{n+1-k}_\xi}\check\psi^k_\xi},
\end{multline}
with integration constant\vspace*{-3pt}
\begin{equation}\label{eq:expec-zeta}
\expecm{\check \zeta^{n,0}_\xi}:=-\sum_{k=1}^{n}\sum_{l=2}^{k+2}(\overline{\check \lambda_\xi^l/\check \lambda_\xi^2})\,\Bexpec{\overline{\check\psi^{k+2-l}_\xi}\check\zeta^{n-k,0}_\xi}.
\end{equation}
With this choice of the constant, the above equation for $\check\zeta^{n,0}_\xi$ is indeed well-posed by the Fredholm alternative as it iteratively ensures that the right-hand side has vanishing average,\vspace*{-3pt}
\begin{equation}\label{eq:chocolate}
\expec{i\xi\cdot\Aa\big(\nabla\check\zeta^{n-1,0}_\xi+i\xi\check\zeta^{n-2,0}_\xi\big)}+\sum_{k=0}^{n+1}\expec{\overline{\check\psi^{n+1-k}_\xi}\check\psi^k_\xi}=0.
\end{equation}
{Factoring out powers of $i\xi$,} we may then define the real-valued symmetric tensor field~$\zeta^{n,0}$ such that\vspace*{-3pt}
\[\zeta^{n,0} \odot (i\xi)^{\otimes (n+1)}:= \check\zeta^{n,0}_{\xi}.\qedhere\]
\end{lem}


\begin{proof}
It suffices to show by induction that~\eqref{eq:chocolate} holds for all $n\ge0$.
For $k\ge0$, testing the defining equation~\eqref{e.spec-1} for $\check\psi^{k+1}_\xi$ with $\check\zeta^{n,0}_\xi$, we get after averaging and integrating by parts,\vspace*{-3pt}
\begin{multline*}
\expec{\overline{\check \psi^{k+1}_\xi} (-\nabla \cdot \Aa \nabla \check\zeta^{n,0}_\xi)}\\[-3pt]
= \expec{\overline{\check\psi^k_\xi} \Big(\nabla \cdot (\Aa i \xi \check\zeta_\xi^{n,0})+i \xi \cdot \Aa\nabla \check\zeta_\xi^{n,0}\Big) +\overline{\check \psi_\xi^{k-1}}(i\xi \cdot\Aa i\xi)\check \zeta_\xi^{n,0}}
+\sum_{l=2}^{k+1}\overline{\check \lambda_\xi^l}\,\Bexpec{\overline{\check \psi_\xi^{k+1-l}}\check\zeta_\xi^{n,0}}.
\end{multline*}
Now inserting the defining equation~\eqref{e.zeta0n} for $\check\zeta_\xi^{n,0}$, and reorganizing the terms, this becomes\vspace*{-3pt}
\begin{multline*}
\expec{\overline{\check \psi_\xi^{k+1}} \Big(\nabla\cdot(\Aa i\xi\check\zeta_\xi^{n-1,0})+i\xi\cdot\Aa(\nabla\check\zeta_\xi^{n-1,0}+i\xi\check\zeta_\xi^{n-2,0})\Big)
+\overline{\check\psi_\xi^k} (i \xi \cdot \Aa i\xi)\check\zeta_\xi^{n-1,0} }\\
= \expec{\overline{\check\psi_\xi^k} \Big(\nabla \cdot (\Aa i \xi \check\zeta_\xi^{n,0})+ i \xi \cdot \Aa(\nabla \check\zeta_\xi^{n,0}+i\xi\check\zeta_\xi^{n-1,0})\Big)
+\overline{\check \psi_\xi^{k-1}}(i\xi \cdot\Aa i\xi)\check \zeta_\xi^{n,0}}\\
+\expec{\overline{\check \psi_\xi^{k+1}}\check\psi_\xi^{n+1}}
+\sum_{l=2}^{k+1}\overline{\check \lambda_\xi^l}\,\Bexpec{ \overline{\check \psi_\xi^{k+1-l}}\check\zeta_\xi^{n,0}}.
\end{multline*}
Iterating this identity $(n+1)$-times, starting from $k=0$, and using that $\psi^0=1$ and that $\expec{\psi^m}=0$ for all $m\geq 1$, we deduce after straightforward simplifications\vspace*{-3pt}
\begin{equation*}
\expec{i\xi \cdot \Aa(\nabla \check\zeta_\xi^{n,0}+i\xi\check\zeta_\xi^{n-1,0})}+\sum_{k=0}^{n+2}\expec{\overline{\check \psi_\xi^{n+2-k}}\check\psi_\xi^{k}}
=-\sum_{k=2}^{n+1}\sum_{l=2}^{k}\overline{\check \lambda_\xi^l}\,\Bexpec{ \overline{\check \psi_\xi^{k-l}}\check\zeta_\xi^{n+1-k,0}},
\end{equation*}
or equivalently,\vspace*{-3pt}
\begin{multline*}
\expec{i\xi \cdot \Aa(\nabla \check\zeta_\xi^{n,0}+i\xi\check\zeta_\xi^{n-1,0})}+\sum_{k=0}^{n+2}\expec{\overline{\check \psi_\xi^{n+2-k}}\check\psi_\xi^{k}}\\[-5pt]
={-}\overline{\check \lambda_\xi^2}\,\expecm{\check\zeta_\xi^{n-1,0}}
{-}\sum_{k=1}^{n-1}\sum_{l=2}^{k+2}\overline{\check \lambda_\xi^l}\,\Bexpec{ \overline{\check \psi_\xi^{k+2-l}}\check\zeta_\xi^{n-k-1,0}}.
\end{multline*}
Now the choice~\eqref{eq:expec-zeta} of the integration constant for $\check\zeta_\xi^{n-1,0}$ precisely ensures that the right-hand side vanishes, which proves that the identity~\eqref{eq:chocolate} holds with $n$ replaced by $n+1$.
\end{proof}

\end{document}
